\chapter{Introducción}
\label{cap:introduccion}
El proyecto se desarrolla en el marco de un trabajo fin de máster dentro del programa de Máster Universitario en Ingeniería de Control, Automatización y Robótica (INCAR) impartido por la Universidad del País Vasco (EHU). El proyecto se desarrolla en la sede de INGEBIT S.L. empresa especializada en ofrecer soluciones en los diferentes ámbitos de la industria 4.0. ubicada en Barakaldo, Bizkaia, España. \par

A inicios de siglo la generación energética en España se basaba principalmente en grandes centrales térmicas y nucleares, con una participación limitada de fuentes renovables. Sin embargo, el auge de la industria 4.0  junto con el creciente problema del cambio climático, ha impulsado que el modelo de generación de energía tradicionalmente monolítico pase a sistemas de generación sostenibles. Estos sistemas son de menor potencia de generación por lo que se requiere un mayor volumen y distribución de ellos.\par

En este contexto, este trabajo fin de máster consiste en el diseño e implementación de un sistema SCADA para la gestión integral de una red de generación de energía sostenible. El sistema permitirá la monitorización centralizada, el control automatizado y la optimización de recursos energéticos, incorporando capacidades de mantenimiento predictivo basadas en análisis de datos en tiempo real.\par

El proyecto se divide en varias fases, incluyendo la planificación, el diseño del sistema SCADA, la implementación y las pruebas finales. Durante el desarrollo  se aplican los conocimientos adquiridos a lo largo del máster, así como las habilidades técnicas y de gestión de proyectos necesarias para llevar a cabo un proyecto de esta envergadura.\par